% GENERATED FILE. Edit canonical lessons, then run npm run generate:books.
% canonical-source-hash: fnv1a64:e9d2812ecea6a07a
% canonical-lessons: TE-C269-dabba, TE-C269-kancam, TE-C269-jalleda, TE-C269-cekka, TE-C269-ittadi

\chapter{Tin, Eating plate, Sieve, Wood, Brass}
\label{ch:te-tin-eating-plate-sieve-wood-brass}
\hlchaptermodality{269}

\begin{chapteropening}
\textbf{By the end of this chapter:} \emph{I can say a tin, an eating plate, a sieve, wood and brass.}

Complete the last lesson of chapter 269: I can say a tin, an eating plate, a sieve, wood and brass.
\end{chapteropening}

\section[\texorpdfstring{\te{డబ్బా} (ḍabbā)}{డబ్బా (ḍabbā)}]{\te{డబ్బా} (ḍabbā) — a tin, a can}

\label{lesson:TE-C269-dabba}



\begin{quote}
Before the new one: say the Telugu for the internet, then the Telugu for a website.
\end{quote}

\subsection*{You'll want to know: \te{డబ్బా}}
\textbf{\te{డబ్బా}} — \emph{ḍabbā} — \textquotedblleft{}a tin, a can\textquotedblright{}.

\begin{grammarlens}[title={What you've built}]
One more word for this chapter.
\end{grammarlens}

\subsection*{Your turn}
\begin{itemize}
\raggedright
  \item \emph{Say it:} \emph{ḍabbā}
  \item \emph{Say it:} \emph{ḍabbā}, once more
  \item \emph{Say it:} say \emph{ḍabbā}
  \item \emph{Recall:} say the Telugu for the internet, then the Telugu for a website, then say \emph{ḍabbā} again
\end{itemize}

\begin{tcolorbox}[breakable,colback=teal!4,colframe=teal!35!black,title={Before you move on}]
What does \te{డబ్బా} mean? (\textquotedblleft{}A tin, a can\textquotedblright{}.) Say it once more.
\end{tcolorbox}
\section[\texorpdfstring{\te{కంచం} (kañcaṁ)}{కంచం (kañcaṁ)}]{\te{కంచం} (kañcaṁ) — an eating plate}

\label{lesson:TE-C269-kancam}



\begin{quote}
Before the new one: say the Telugu for a website, then the Telugu for a tin.
\end{quote}

\subsection*{You'll want to know: \te{కంచం}}
\textbf{\te{కంచం}} — \emph{kañcaṁ} — \textquotedblleft{}an eating plate\textquotedblright{}.

\begin{grammarlens}[title={What you've built}]
One more word for this chapter.
\end{grammarlens}

\subsection*{Your turn}
\begin{itemize}
\raggedright
  \item \emph{Say it:} \emph{kañcaṁ}
  \item \emph{Say it:} \emph{kañcaṁ}, once more
  \item \emph{Say it:} say \emph{kañcaṁ}
  \item \emph{Recall:} say the Telugu for a website, then the Telugu for a tin, then say \emph{kañcaṁ} again
\end{itemize}

\begin{tcolorbox}[breakable,colback=teal!4,colframe=teal!35!black,title={Before you move on}]
What does \te{కంచం} mean? (\textquotedblleft{}An eating plate\textquotedblright{}.) Say it once more.
\end{tcolorbox}
\section[\texorpdfstring{\te{జల్లెడ} (jalleḍa)}{జల్లెడ (jalleḍa)}]{\te{జల్లెడ} (jalleḍa) — a sieve}

\label{lesson:TE-C269-jalleda}



\begin{quote}
Before the new one: say the Telugu for a tin, then the Telugu for an eating plate.
\end{quote}

\subsection*{You'll want to know: \te{జల్లెడ}}
\textbf{\te{జల్లెడ}} — \emph{jalleḍa} — \textquotedblleft{}a sieve\textquotedblright{}.

\begin{grammarlens}[title={What you've built}]
One more word for this chapter.
\end{grammarlens}

\subsection*{Your turn}
\begin{itemize}
\raggedright
  \item \emph{Say it:} \emph{jalleḍa}
  \item \emph{Say it:} \emph{jalleḍa}, once more
  \item \emph{Say it:} say \emph{jalleḍa}
  \item \emph{Recall:} say the Telugu for a tin, then the Telugu for an eating plate, then say \emph{jalleḍa} again
\end{itemize}

\begin{tcolorbox}[breakable,colback=teal!4,colframe=teal!35!black,title={Before you move on}]
What does \te{జల్లెడ} mean? (\textquotedblleft{}A sieve\textquotedblright{}.) Say it once more.
\end{tcolorbox}
\section[\texorpdfstring{\te{చెక్క} (cekka)}{చెక్క (cekka)}]{\te{చెక్క} (cekka) — wood, a plank}

\label{lesson:TE-C269-cekka}



\begin{quote}
Before the new one: say the Telugu for an eating plate, then the Telugu for a sieve.
\end{quote}

\subsection*{You'll want to know: \te{చెక్క}}
\textbf{\te{చెక్క}} — \emph{cekka} — \textquotedblleft{}wood, a plank\textquotedblright{}.

\begin{grammarlens}[title={What you've built}]
One more word for this chapter.
\end{grammarlens}

\subsection*{Your turn}
\begin{itemize}
\raggedright
  \item \emph{Say it:} \emph{cekka}
  \item \emph{Say it:} \emph{cekka}, once more
  \item \emph{Say it:} say \emph{cekka}
  \item \emph{Recall:} say the Telugu for an eating plate, then the Telugu for a sieve, then say \emph{cekka} again
\end{itemize}

\begin{tcolorbox}[breakable,colback=teal!4,colframe=teal!35!black,title={Before you move on}]
What does \te{చెక్క} mean? (\textquotedblleft{}Wood, a plank\textquotedblright{}.) Say it once more.
\end{tcolorbox}
\section[\texorpdfstring{\te{ఇత్తడి} (ittaḍi)}{ఇత్తడి (ittaḍi)}]{\te{ఇత్తడి} (ittaḍi) — brass}

\label{lesson:TE-C269-ittadi}



\begin{quote}
Before the new one: say the Telugu for a sieve, then the Telugu for wood.
\end{quote}

\subsection*{You'll want to know: \te{ఇత్తడి}}
\textbf{\te{ఇత్తడి}} — \emph{ittaḍi} — \textquotedblleft{}brass\textquotedblright{}.

\begin{grammarlens}[title={What you've built}]
One more word for this chapter.
\end{grammarlens}

\subsection*{Your turn}
\begin{itemize}
\raggedright
  \item \emph{Say it:} \emph{ittaḍi}
  \item \emph{Say it:} \emph{ittaḍi}, once more
  \item \emph{Say it:} say \emph{ittaḍi}
  \item \emph{Recall:} say the Telugu for a sieve, then the Telugu for wood, then say \emph{ittaḍi} again
\end{itemize}

\begin{tcolorbox}[breakable,colback=teal!4,colframe=teal!35!black,title={Before you move on}]
What does \te{ఇత్తడి} mean? (\textquotedblleft{}Brass\textquotedblright{}.) Say it once more.
\end{tcolorbox}
